% Build with: xelatex -interaction=nonstopmode -halt-on-error -output-directory=assets assets/CV_source.tex
\documentclass[letterpaper,10pt]{article}

\usepackage[T1]{fontenc}
\usepackage[utf8]{inputenc}
\usepackage{mathpazo}
\usepackage[margin=0.65in]{geometry}
\usepackage[hidelinks]{hyperref}
\usepackage{tabularx}
\usepackage{array}
\usepackage{enumitem}
\usepackage{titlesec}
\usepackage{fancyhdr}
\usepackage{xcolor}
\usepackage[none]{hyphenat}
\usepackage{microtype}

\pagestyle{fancy}
\fancyhf{}
\renewcommand{\headrulewidth}{0pt}
\fancyfoot[L]{\footnotesize Updated September 2026}
\fancyfoot[C]{\footnotesize Md Ashikul Haque \textbullet\ Curriculum Vitae \textbullet\ \thepage}

\setlength{\parindent}{0pt}
\setlength{\parskip}{0pt}
\setlength{\tabcolsep}{0pt}
\emergencystretch=1.5em

\titleformat{\section}
  {\large\bfseries\scshape}
  {}
  {0pt}
  {}
  [\vspace{-0.45em}\rule{\linewidth}{0.4pt}\vspace{0.1em}]
\titlespacing*{\section}{0pt}{0.72em}{0.28em}

\setlist[itemize]{leftmargin=1.2em,itemsep=1pt,topsep=2pt,parsep=0pt,partopsep=0pt,label=\textbullet}
\setlist[enumerate]{leftmargin=1.8em,itemsep=2pt,topsep=2pt,parsep=0pt,partopsep=0pt}

\newcolumntype{Y}{>{\raggedright\arraybackslash}X}
\newcommand{\entry}[2]{\begin{tabularx}{\linewidth}{@{}Y r@{}}#1 & #2\end{tabularx}}
\newcommand{\subentry}[1]{#1\par}
\newcommand{\pubitem}[2]{\noindent\makebox[1.8em][l]{[#1]}#2\par}

\begin{document}\sloppy

\begin{center}
{\LARGE\bfseries MD ASHIKUL HAQUE}\par
\vspace{4pt}
{\small
\href{mailto:mdashikul.haque@utdallas.edu}{mdashikul.haque@utdallas.edu}
\enspace | \enspace
+1 (515) 598 6458
\enspace | \enspace
\href{https://ashikul-haque.github.io}{ashikul-haque.github.io}\par
Google Scholar:
\href{https://scholar.google.com/citations?user=11vTHPYAAAAJ}{scholar.google.com/citations?user=11vTHPYAAAAJ}
}
\end{center}

\section{Education}

\entry{\textbf{PhD, Computer Science}}{Expected May 2027}\par
University of Texas at Dallas, USA\par
Dissertation: \emph{Dependable Low-Power Wide-Area Networking under Contention: Jamming, Coexistence, and Deadlines}\par
Advisor: Prof. Abusayeed Saifullah

\vspace{4pt}

\entry{\textbf{MSc, Computer Science}}{2024}\par
Wayne State University, USA

\vspace{4pt}

\entry{\textbf{BSc, Computer Science and Engineering}}{2018}\par
Bangladesh University of Engineering and Technology, Bangladesh

\section{Research Interests}

Anti-Jamming and Physical-Layer Security for LPWANs \textbullet\ Coexistence among Learning-Enabled Networks \textbullet\ Real-Time and Mixed-Criticality Scheduling \textbullet\ Cross-Technology and Direct-to-Satellite Links \textbullet\ Deadline-Aware Scheduling for LLM Agents and Connected Vehicles

\vspace{4pt}

\textbf{Research Summary.} I build low-power wireless systems that stay correct and on time when the resource they depend on is contested by an attacker, by neighboring networks, or by a deadline. The end device is the most constrained part of an Internet of Things (IoT) system, so the defense, the coordination, and the scheduling belong in the infrastructure. My work covers physical-layer defenses that keep LoRa networks decoding under jamming attacks, learning-based coexistence management for dense low-power wide-area network (LPWAN) deployments, medium access control (MAC) and mixed-criticality scheduling for industrial wireless networks, cross-technology communication from Long Range Frequency Hopping Spread Spectrum (LR-FHSS) to LoRa, and deadline-aware token scheduling for large language model (LLM) agents that drive physical actuators. I validate these designs on software-defined radio (SDR) prototypes with GNU Radio and USRP B200, on commercial off-the-shelf (COTS) LoRa hardware and multi-node testbeds, and in PyTorch, NS-3, and CARLA at scale.

\section{Conference Publications}

First author of every paper except [2], for which I am the second author.

\vspace{3pt}

\pubitem{9}{\textbf{Md Ashikul Haque}, Abusayeed Saifullah. \emph{Deadline-Aware Token Scheduling for Physical-I/O LLM Agents}. In Proc. IEEE Real-Time Systems Symposium (RTSS 2026). Acceptance rate: 14\%.}
\vspace{2pt}
\pubitem{8}{\textbf{Md Ashikul Haque}, Aakriti Jain, Venkata Prashant Modekurthy, Abusayeed Saifullah. \emph{Enabling Cross Technology Communication from LR-FHSS to LoRa}. In Proc. ACM/IEEE International Conference on Embedded Artificial Intelligence and Sensing Systems (SenSys 2026), CPS-IoT Week 2026. Acceptance rate: 19.8\% (52 of 263).}
\vspace{2pt}
\pubitem{7}{\textbf{Md Ashikul Haque}, Abusayeed Saifullah. \emph{Decoding LoRa Packets under Collaborative Jamming Attacks}. In Proc. 23rd International Conference on Embedded Wireless Systems and Networks (EWSN 2026). Acceptance rate: 23.5\% (8 of 34).}
\vspace{2pt}
\pubitem{6}{\textbf{Md Ashikul Haque}, Abusayeed Saifullah. \emph{ReMix: Resource-Reclaiming Mixed-Criticality Scheduling for Industrial IoT}. In Proc. The 35th International Conference on Computer Communications and Networks (ICCCN 2026). Acceptance rate: 25\%.}
\vspace{2pt}
\pubitem{5}{\textbf{Md Ashikul Haque}, Abusayeed Saifullah. \emph{Mitigating Jamming Attacks in LoRa Networks: A Defense Strategy against LoRa-Based Jammers}. In Proc. ACM International Symposium on Theory, Algorithmic Foundations, and Protocol Design for Mobile Networks and Mobile Computing (MobiHoc 2025). Acceptance rate: 23\% (39 of 169).}
\vspace{2pt}
\pubitem{4}{\textbf{Md Ashikul Haque}, Abusayeed Saifullah, Haibo Zhang. \emph{Deep Reinforcement Learning Based Coexistence Management in LPWAN}. In Proc. IEEE Conference on Computer Communications (INFOCOM 2025). Acceptance rate: 18.6\% (272 of 1458).}
\vspace{2pt}
\pubitem{3}{\textbf{Md Ashikul Haque}, Abusayeed Saifullah. \emph{Handling Jamming Attacks in a LoRa Network}. In Proc. ACM/IEEE Conference on Internet of Things Design and Implementation (IoTDI 2024), CPS-IoT Week 2024.}
\vspace{2pt}
\pubitem{2}{Aakriti Jain, \textbf{Md Ashikul Haque}, Abusayeed Saifullah, Haibo Zhang. \emph{Burst-MAC: A MAC Protocol for Handling Burst Traffic in LoRa Network}. In Proc. IEEE Real-Time Systems Symposium (RTSS 2024).}
\vspace{2pt}
\pubitem{1}{\textbf{Md Ashikul Haque}, Abusayeed Saifullah. \emph{A Game-Theoretic Approach for Mitigating Jamming Attacks in LPWAN}. In Proc. International Conference on Embedded Wireless Systems and Networks (EWSN 2023). To our knowledge, the first anti-jamming work in LPWAN.}

\section{Workshop Publications}

\begin{enumerate}[label={[\arabic*]}]
\item \textbf{Md Ashikul Haque}, Abusayeed Saifullah, Haibo Zhang. \emph{Coordinating Coexisting Learning Agents in Shared Spectrum via Parameter Space Complementarity}. Agents in the Wild: Safety, Security, and Beyond (ICLR 2026).
\item \textbf{Md Ashikul Haque}, Dali Ismail, Abusayeed Saifullah. \emph{Scalable Real-Time Control in Industrial Cyber-Physical Systems}. In Proc. Workshop on Software-Defined Networking for Cyber-Physical System Automation (SDN-CPS 2024), co-located with ICDCN 2024, pp. 353--358.
\end{enumerate}

% Manuscripts Under Review is deliberately kept out of the public CV.
% The section lives in assets/CV_private.tex, which git does not track.
% Build the private version with:
%   xelatex -interaction=nonstopmode -output-directory=assets assets/CV_private.tex
\ifdefined\PrivateExtras\PrivateExtras\fi

\section{Selected Talks}

\entry{\emph{ReMix: Resource-Reclaiming Mixed-Criticality Scheduling for Industrial IoT}}{Jul 2026}\par
IEEE ICCCN 2026, Honolulu, Hawaii, USA

\vspace{4pt}

\entry{\emph{Enabling Cross Technology Communication from LR-FHSS to LoRa}}{May 2026}\par
ACM/IEEE SenSys 2026, CPS-IoT Week 2026

\vspace{4pt}

\entry{\emph{Mitigating Jamming Attacks in LoRa Networks: A Defense Strategy against LoRa-Based Jammers}}{Oct 2025}\par
ACM MobiHoc 2025, Rice University, Houston, USA

\vspace{4pt}

\entry{\emph{Handling Jamming Attacks in a LoRa Network}}{May 2024}\par
ACM/IEEE IoTDI 2024, Hong Kong, China

\section{Research and Professional Experience}

\entry{\textbf{Graduate Research Assistant, Dept. of Computer Science}}{2022 to 2026}\par
Wayne State University, USA (Fall 2022 to Summer 2025) and The University of Texas at Dallas, USA (Fall 2025 and Summer 2026) | Advisor: Prof. Abusayeed Saifullah
\begin{itemize}
\item Published 9 conference papers and 2 workshop papers in wireless networking and real-time systems, 10 of them as first author.
\item Designed gateway-side and network-server-side defenses that recover LoRa packets under single, LoRa-based, collaborative, and spatially distributed jammers, with outdoor evaluation on USRP B200 gateways and commodity LoRa nodes.
\item Built learning-based coexistence management at the LoRa Network Server and disclosure-free coordination among LoRaWAN operators, evaluated in PyTorch, NS-3, and a demodulator-level simulation of 32,000 devices across 8 operators.
\item Designed real-time medium access and mixed-criticality scheduling for LoRa and industrial wireless sensor-actuator networks, evaluated on LoRa devices and a 15-node Tmote Sky testbed, and extended the same scheduling method to LLM serving on vLLM/ROCm and to cooperative perception on a three-car MentorPi testbed and CARLA.
\item Built cross-technology communication from LR-FHSS to LoRa, per-message delivery certificates, and signature-bound frame recovery for direct-to-satellite LR-FHSS on GNU Radio and USRP.
\end{itemize}

\vspace{4pt}

\entry{\textbf{Assistant Engineer, ICT Division}}{2020 to 2021}\par
Dhaka Electric Supply Company, Bangladesh
\begin{itemize}
\item Maintained ICT infrastructure with monitoring, incident response, and reliability focused operations for distributed utility systems.
\end{itemize}

\vspace{4pt}

\entry{\textbf{Software Engineer, IoT Division}}{2019 to 2020}\par
Samsung R\&D Institute Bangladesh, Bangladesh
\begin{itemize}
\item Worked across both the user interface and the cloud integration layers of the SmartThings iOS application, which Samsung reported serving more than 45 million monthly active users and 5,000 supported device types in October 2019.
\item Developed production-grade C++ components in the Samsung Cloud Client (SCClient), the shared library that handles cloud connectivity and device communication for SmartThings, and exposed them to iOS through a Swift wrapper.
\item Built the device onboarding dashboard through which users register, authenticate, and configure new devices.
\item Automated the entire build, packaging, and release workflow of the SCClient iOS team in Bash, which removed the manual steps from every release and made each build reproducible.
\end{itemize}

\section{Teaching}

\entry{\textbf{Graduate Teaching Assistant, Dept. of Computer Science}}{Spring and Fall 2026}\par
The University of Texas at Dallas, USA\par
Courses: CS 5390 Computer Networks; CS 6352 Performance of Computer Systems and Networks
\begin{itemize}
\item Graded assignments, exams, and project deliverables for both courses.
\item Helped students understand course concepts through office hours, question answering, and discussion of networking and systems topics.
\item Supported course projects by clarifying requirements, troubleshooting implementations, and providing feedback on design and debugging.
\end{itemize}

\vspace{4pt}

\entry{\textbf{Guest Lecturer, Dept. of Computer Science}}{2025}\par
The University of Texas at Dallas, USA\par
Course: Internet of Things (graduate). Delivered multiple lectures on LoRa networking and network simulation in NS-3.

\vspace{4pt}

\entry{\textbf{Lecturer, Dept. of Computer Science}}{2018 to 2019}\par
Millennium University, Bangladesh\par
Courses: Structured Programming; Data Communication. Instructor of record for both courses: designed lectures, laboratory assignments, and examinations.

\section{Mentoring}

\entry{\textbf{Nasif Ahmed, PhD student, The University of Texas at Dallas}}{January 2025 to present}\par
Junior doctoral student in our laboratory. I mentor him on identifying new problems and approaching their solutions, on experiments with USRP radios and commercial off-the-shelf (COTS) LoRa devices, and on writing systems papers.

\section{Awards and Honors}

\entry{\textbf{Dean's List Award, Wayne State University}}{2024}
\entry{\textbf{Professional Travel Award, College of Engineering, Wayne State University}}{2024}
\entry{\textbf{ACM SIGBED Travel Grant, CPS-IoT Week, San Antonio, Texas, USA}}{2023}

\section{Research Funding and Proposals}

\begin{itemize}
\item Contributed to the writing of an NSF proposal on real-time cooperative perception for connected vehicles, with Prof. Abusayeed Saifullah as principal investigator \hfill 2026
\item My doctoral research has been supported by five NSF awards (CNS-2301757, CAREER-2306486, CNS-2306745, CNS-2601685, and CNS-2602744) and one ONR award (N00014-23-1-2151), all with Prof. Abusayeed Saifullah as principal investigator \hfill 2022 to present
\end{itemize}

\section{Professional Service}

\textbf{Conference Referee}
\begin{itemize}
\item ACM/IEEE Intl. Conf. on Embedded Artificial Intelligence and Sensing Systems (SenSys) \hfill 2022, 2024, 2025, 2026
\item IEEE Real-Time Systems Symposium (RTSS) \hfill 2023, 2024, 2025
\item IEEE Real-Time and Embedded Technology and Applications Symposium (RTAS) \hfill 2023, 2024
\item Intl. Conf. on Embedded Wireless Systems and Networks (EWSN) \hfill 2023, 2024, 2025
\item IEEE Intl. Conf. on Distributed Computing in Smart Systems and the IoT (DCOSS-IoT) \hfill 2024, 2026
\item IEEE Intl. Conf. on Embedded and Real-Time Computing Systems and Applications (RTCSA) \hfill 2024, 2025
\item IEEE Intl. Conf. on Sensing, Communication, and Networking (SECON) \hfill 2026
\end{itemize}

\textbf{Journal Referee}
\begin{itemize}
\item IEEE/ACM Transactions on Networking \hfill 2024, 2025
\item Pervasive and Mobile Computing (Elsevier) \hfill 2023, 2024
\end{itemize}

\textbf{Other Service}
\begin{itemize}
\item Student Volunteer, CPS-IoT Week 2023, San Antonio, Texas, USA \hfill 2023
\end{itemize}

\section{Technical Skills}

\begin{tabularx}{\linewidth}{@{}l@{\hspace{1em}}Y@{}}
\textbf{Programming:} & C, C++, Python, MATLAB, SQL, Swift, Bash \\
\textbf{Wireless and SDR:} & GNU Radio, USRP B200, LoRa and LR-FHSS physical layers, LoRaWAN, COTS LoRa nodes and gateways, Tmote Sky \\
\textbf{Learning and simulation:} & PyTorch, Pandas, NS-3, CARLA, vLLM \\
\textbf{Systems and tools:} & Linux, Docker, Git, LaTeX \\
\end{tabularx}

\section{References}

\begin{tabularx}{\linewidth}{@{}X X@{}}
\textbf{Dr. Abusayeed Saifullah} (advisor) & \textbf{Dr. Haibo Zhang} \\
Associate Professor, Dept. of Computer Science & Associate Professor, School of Computer Science \\
University of Texas at Dallas & and Engineering, UNSW Sydney \\
\href{mailto:abusayeed.saifullah@utdallas.edu}{abusayeed.saifullah@utdallas.edu} & \href{mailto:haibo.zhang@unsw.edu.au}{haibo.zhang@unsw.edu.au} \\
+1 (972) 883.4171 & \\
\end{tabularx}

\vspace{8pt}

\begin{tabularx}{\linewidth}{@{}X@{}}
\textbf{Dr. Venkata Prashant Modekurthy} \\
Assistant Professor, Dept. of Computer Science \\
University of Nevada, Las Vegas (UNLV) \\
\href{mailto:prashant.modekurthy@unlv.edu}{prashant.modekurthy@unlv.edu} \\
+1 (702) 774.3405 \\
\end{tabularx}

\end{document}
